Nog niet lang geleden
leefden tussen steden

veroordeelde lieden
in ruige gebieden.

Ooit vogelvrij verklaard,
zijn ze van huis en haard,

uit hun gekend gehucht,
met tegenzin gevlucht.

Lijfeigene blijven,
was niet te beschrijven

qua onmenselijk leed.
Dus wat menigeen deed,

was eraan ontsnappen
door plotsklaps te kappen

met de harde knechtschap.
Deze louche misstap

leidde tot verbanning
en tot een verklaring

dat de stoute knechten
hun geringe rechten

gans hadden verkwanseld.
Catastrofaal afgeranseld

werd de vluchter bij vangst.
Derhalve uit doodsangst

verscholen velen zich.
Aldoende zwierven tig

vogelvrijverklaarden
als het slijk der aarde

door nabije wouden
om pas halt te houden

in de veengebieden.
Daar konden de lieden

zich soms veilig achten.
Dat viel te betrachten.

Al was de natuur mooi,
ook daar viel men ten prooi.

Geen wet die je beschermt,
geen macht die zich ontfermt

over jou en je vrienden.
Sommigen verdienden

aan de vluchters een duit.
Trokken er vaak op uit

de poet te vangen.
Vanuit het verlangen

naar een voorbeeld stellen,
ging men hen bestellen.

Zo werden struikrovers
legale uitslovers

en geen mensenstropers.
De angst van meelopers

voor de libertijnen
deed hen snel verdwijnen,

onder druk van Pascal,
tegen moreel verval.

Gedoogd door aanklagers
konden premiejagers

vogelvrijen strikken
en geheel beschikken

over de eigen poet.
Meestal werd het goed

gegeven aan de klant.
Die maakte ze van kant

in een publieke show.
Tot zover dit intro.

Laat dit essay gaan naar
het bijzondere jaar

1711.
Een Duitser hielp zichzelf

in zijn jagerskleding
bij zijn jagerswoning,

vanwaar hij rondzwierf
en zijn faam verwierf

met de vele prooien
die hij wist te kooien.

Hij mag geen naam hebben!
Maar wel zal ik reppen

over zijn mensenjacht.
Doch de ware toedracht,

is ook mij onbekend.
Dat bent u vast gewend.

Nooit in zijn eigen stad
bracht hij zijn verse schat.

Tenminste zo meenden
alle entertainden

bij de terechtstelling
van een vluchteling.

Elders bracht hij zijn buit.
Daar betaalde men hem uit.

Hij trok daar de aandacht.
Doch kwam er geen aanklacht

toen alerte leden
uit nabije steden

melden aan iedereen,
dat hij zich alleen

voor geknepen mannen
echt leek in te spannen.

Waarvoor men hem roemde.
De critici noemden

de “verdwijntovenaar”
echter een “tevenaar”.

Voor juffrouwen wellicht
kneep hij een oogje dicht.

Als hofleverancier
van de regio alhier

betekende dat wat.
Een gevlucht meisje had

ter plekke in wezen
heel weinig te vrezen.

Dus gingen ze massaal
met hun lot aan de haal.

Al leek de illusie,
door angst in remissie,

als sluier te vallen,
bleek het zich te stallen

in het onbewuste,
waar het even ruste.

Bij controleverlies
herhaalden de ficties,

die Maya verspreidde
en de mens verleidden.

Angst, slecht als raadgever,
kon als richtinggever

dit ego verraden
in al zijn façaden,

maar was even spoorloos.
Vandaar dat een vrouw koos

voor teveel vertrouwen
en moest dat berouwen.

Hij dronk meer dan gangbaar
in een tijd waar schijnbaar

water alleen gekookt
en na gisten gestookt

kon worden geschonken.
Hij was dus vaak dronken.

De bejaarde zuiplap
was een lelijk heerschap.

Hij verloor ooit een oog
door een pijl uit een boog

van een van zijn prooien.
Nu een van zijn dooien.

Buiten een slecht gebit
was hij verder topfit.

Toen hij zijn stad verliet,
richting heidegebied,

vergezelde nog lang
zijn kalme wandelgang

de geuren van taarten.
Het was daar Sint Maarten.

Kinderen in optocht
markeerden zijn aftocht

in vol temperament.
Wachtend op het moment

nadat het morgenrood
met water in de sloot

overgaat in avondgoud.
Volgend werd dan gesjouwd

achter lampionnen,
die hun reis begonnen

op jacht naar lekkernij.
Een moeder leek niet blij

met de voorbereiding
toen een kaars niet goed hing.

“Wat ben je toch een oen!
Laat mij het maar zelf doen.

Straks krijgen we nog brand”
De toorts kwam in haar hand

en het schoof met een zwier
naar het gesmeerd papier

rondom een lampion,
die hing aan het plafond

met een brandende kaars.
Het gaf de knutselaars

een toepasselijk licht.
Met een onthutst gezicht
keek de kleine madam.
De kaarslamp vatte vlam

en stortte neer, brandend
op de tafel landend

met alle lampions.
Je hoorde binnensmonds

moeder zich verdommen.
Het kind bleef verstommen.

Hij miste deze pech.
Hij baande zich zijn weg

al met zijn wandelstok.
Een dichte landmist trok

over het vredige veen.
Het vergezicht verdween

en werd “misterieus”.
De mens werd curieus.

Een waterkoude bries,
die onbekommerd blies,

verspreidde de kilte
en verbrak de stilte.

Dit weer guur, grijs en grauw,
met die snijdende kou,

meden premiejagers
als verwende klagers.

Patserige rijkdom
verworven met weedom

had hun flinkheid geroofd.
Comfortabel verdoofd.

De jager bleef niet thuis.
Zo arm als een kerkmuis

was onze hoofdpersoon.
Hij gaf zijn vindersloon

aan elke herbergier
voor een paar liters bier.

Voor hem was dit weertype
geen grond om te piepen.

Immers vogelvrijen
konden zich verblijden

met de klare gedachte
dat men mocht verwachten

dat thans hun belagers,
de schadelijke jagers,

met barokke lijven,
liever binnen blijven.

Afwezig leek Phobos
van elke fobie los.

Dus werden ze slordig
en overijld lichtvaardig.

Dat bood schakers kansen.
Hij ging zich verschansen.

Welke prooi vrij van angst
werd zijn volgende vangst?
